\documentclass[aps,prb,onecolumn,nofootinbib,superscriptaddress]{revtex4-2}
\usepackage{amsmath,amssymb,amsthm,mathtools,bm}
\usepackage{booktabs}
\usepackage[colorlinks=true,linkcolor=blue,citecolor=blue,urlcolor=blue]{hyperref}
\usepackage{microtype}

\newcommand{\Tr}{\operatorname{Tr}}
\newcommand{\I}{\mathbf{1}}
\newcommand{\ket}[1]{\lvert #1\rangle}
\newcommand{\bra}[1]{\langle #1\rvert}
\newcommand{\pathref}[1]{\path{#1}}

\begin{document}

\title{Worker 13 Note: Regularized Choi Transfer Evidence}
\author{Repository synthesis worker 13}
\date{\today}

\begin{abstract}
I inspect the regularized Choi-transfer materials in
\path{colab_regularized_choi_transfer_matrix/}, excluding all paths requested
by the hierarchy.  The evidence supports a coherent operator-level diagnostic:
realized adaptive trajectories are still propagated by the canonical GPU Markov
engine, while an auxiliary cumulative Gaussian Choi covariance is updated by a
regulated rank-one contraction and queried for a finite-sector particle-transfer
gap.  The strongest result is not a long-time physical trajectory gap, but a
controlled short-to-intermediate-time transfer/stability diagnostic with explicit
censoring when the Choi chart becomes numerically unreliable.
\end{abstract}

\maketitle
\setcounter{tocdepth}{2}
\tableofcontents

\section{Source Scope}

I read \path{REPO_POLICY.md} before inspection.  I inspected only lightweight
text, metadata, notebook JSON text, source code, and CSV summaries under
\path{colab_regularized_choi_transfer_matrix/}.  I did not execute notebooks,
edit files, or load binary arrays.  I excluded paths containing
\texttt{haining}, \texttt{haoyu}, or \texttt{erroneous}, and did not read under
\path{repo_synthesis/main_pass/worker_notes}, \path{supervisor_notes}, or
\path{boss}.

\section{Findings}

The theoretical note distinguishes the cumulative operator tensor from the
physical state covariance.  The physical covariance \(G\) is the production
trajectory variable, while the two-leg Choi covariance \(\Sigma\) represents the
realized circuit operator \(\hat C_t\), with \(L,R\) denoting Choi legs rather
than spatial regions.  This distinction is explicit in the derivation and is
essential for interpreting the evidence.

The local branch tensor has the singular rank-one form
\[
K(\eta_1,\eta_2)=
\begin{pmatrix}
\eta_1 P & Q\\
Q & \eta_2 P
\end{pmatrix},
\qquad
P=\ket{\chi}\bra{\chi},\quad Q=\I-P ,
\]
so the Choi contraction requires regulation on \(Q\).  The note introduces
\(K_{LL}^{(\epsilon)}=\eta_1P+\epsilon Q\), takes the Schur-complement limit,
and obtains a finite update when
\[
d=\Tr(P\Sigma_{RR})-\eta_1\neq 0.
\]
The basis-independent update is
\[
\Sigma'=
\begin{pmatrix}
\Sigma_{LL}-\dfrac{\Sigma_{LR}P\Sigma_{RL}}{d}
&
\Sigma_{LR}Q-\dfrac{\Sigma_{LR}P\Sigma_{RR}Q}{d}
\\[8pt]
Q\Sigma_{RL}-\dfrac{Q\Sigma_{RR}P\Sigma_{RL}}{d}
&
\eta_2P+Q\Sigma_{RR}Q-\dfrac{Q\Sigma_{RR}P\Sigma_{RR}Q}{d}
\end{pmatrix}.
\]
The same note proves Hermiticity and pure-Choi involution preservation on
nonsingular branches.

The GPU implementation follows the same structure.  It initializes
\(\Sigma_{LL}=0\), \(\Sigma_{LR}=\I\), \(\Sigma_{RR}=0\), records a singular
tolerance, computes \(d=\bra{\chi}\Sigma_{RR}\ket{\chi}-\eta_1\), records the
minimum denominator context, and either raises or records a
\texttt{denominator\_regularized} event when \(|d|\) is below tolerance.  The
rank-one update is assembled through a resolvent action, matching the
regularized formulation.

The transfer observable is explicitly a finite-sector particle-transfer
diagnostic,
\[
T_p=-(\I+\Sigma_{LL})^{-1}\Sigma_{LR},
\qquad
\Delta_\lambda(t)=\min_j |\log s_j(T_p)/t|.
\]
The helper extracts near-gap modes using a restarted block-Davidson calculation
on \(\Sigma_{LL}^2\), with exact eigensolves used for final spectra or fallback.
Endpoint sectors at \(\Sigma_{LL}=\pm1\) are treated as exact particle
zeros/poles rather than ordinary finite exponents.

\section{Interpretation}

The evidence supports the following reading.  The regularized Choi transfer
object is an operator-history chart layered on top of the physical Markov
trajectory, not a replacement for the physical trajectory covariance.  It is
useful because the cumulative Choi covariance turns products of local adaptive
Kraus events into a transfer spectrum whose near-zero Lyapunov exponents test
whether the realized operator has stable finite-sector transmission channels.

The domain-wall slab runs are the physically targeted transfer/stability
diagnostic.  The homogeneous \(DW=0\) full-basis run is best interpreted as a
control: it uses the same Choi machinery and canonical dynamics path but lacks
the domain-wall slab truncation.  The campaign metadata confirms that the
physical \(G\) evolution continues even after the Choi observable is censored;
therefore the late-cycle disappearance of stable Choi samples is not evidence
that the physical Markov dynamics terminates.  It is evidence that the exact
cumulative Choi chart, under the chosen tolerances and finite-sector extraction,
has become numerically unreliable.

The local CSV summary shows all three available runs retain full Choi stability
for early cycles and then lose all stable samples around cycles \(15\)--\(16\).
For \(N_x=20,N_y=30\), the \(DW=1\) slab gap is small, roughly
\(10^{-2}\)--\(2\times10^{-2}\) over the stable window, while the \(DW=0\) full
control has much larger transient mean gaps after the first few cycles.  For
\(N_y=40\), the \(DW=1\) slab gap is again around \(10^{-2}\) before censoring.
This is consistent with a domain-wall transfer channel remaining closer to
marginality than the homogeneous full-basis control, but only within the
surviving Choi-observable window.

\section{Evidence Table}

\begin{table}[h]
\caption{\label{tab:evidence}Lightweight sources used for the synthesis.}
\begin{ruledtabular}
\begin{tabular}{p{0.34\linewidth}p{0.16\linewidth}p{0.42\linewidth}}
Source & Lines & Evidence \\
\colrule
\pathref{docs/regularized_choi_covariance_contractions.tex} & 63--103 &
Defines the cumulative Choi operator viewpoint and distinguishes it from the
physical state covariance propagated by the canonical engine.\\
\pathref{docs/regularized_choi_covariance_contractions.tex} & 164--246,
455--551 & Derives the singular contraction, regulator, nonsingular condition
\(d\neq0\), and full basis-independent update.\\
\pathref{docs/regularized_choi_covariance_contractions.tex} & 749--775 &
Rewrites the update as regularized resolvent solves, giving the numerical form
used by the implementation.\\
\pathref{docs/regularized_choi_covariance_contractions.tex} & 800--938 &
States Hermiticity, singular-branch limitation, and preservation of the pure
Choi involution.\\
\pathref{src/classA_U1FGTN_gpu.py} & 754--774 &
Initializes cumulative Choi blocks and records singular tolerance and diagnostic
state.\\
\pathref{src/classA_U1FGTN_gpu.py} & 840--916 &
Computes \(d\), detects small denominators, applies resolvent actions, and
updates \(LL,LR,RR\) blocks.\\
\pathref{src/choi_transfer_observables_gpu.py} & 11--18,
43--63 & Records helper metadata, transfer definition, gap definition, endpoint
zero/pole handling, and extraction method.\\
\pathref{src/choi_transfer_observables_gpu.py} & 360--435,
477--635 & Defines the particle Choi gap observer, censoring behavior, retained
diagnostics, and final exact-spectrum path.\\
\pathref{analysis_outputs/gap_vs_cycle/gap_vs_cycle_summary.csv} & 62--76,
122--136 & Shows \(DW=1\) slab mean gaps near \(10^{-2}\) before Choi stability
collapses by cycles \(15\)--\(16\).\\
\pathref{analysis_outputs/gap_vs_cycle/gap_vs_cycle_summary.csv} & 2--17 &
Shows the \(DW=0\) full control has larger transient gaps and also loses all
stable Choi samples by cycle \(16\).\\
\pathref{gpu_data/.../N20x30_DW1.../run_summary.json} & 671--672,
1427--1499, 1716--1778 & Confirms canonical dynamics entry point, censoring
rule, Choi failure mode, slab effectiveness, stable fraction list, and zero
stable final samples.\\
\pathref{gpu_data/.../N20x40_DW1.../run_summary.json} & 891--892,
1647--1719, 1936--2018 & Confirms the same metadata for the \(N_y=40\) slab run.\\
\pathref{gpu_data/.../N20x30_DW0.../run_summary.json} & 1211--1212,
1985--2057, 2292--2354 & Confirms the homogeneous full-basis control,
censoring, and zero stable final samples.\\
\end{tabular}
\end{ruledtabular}
\end{table}

\section{Caveats}

The locally present completed campaign summaries cover \(N_y=30\) for
\(DW=0\), and \(N_y=30,40\) for \(DW=1\).  The notebook text describes a larger
\(N_y=30,40,50\) plan, but I did not treat absent local summaries as completed
evidence.

No full Choi matrices or transfer matrices were inspected; campaign metadata
says they were not saved.  This note therefore relies on documented formulas,
implementation text, JSON metadata, and aggregate CSV summaries.

The late-cycle stable fraction reaches zero in all available runs.  Hence the
final-cycle gap values at \(T=2N_y\) are not physical gap estimates from stable
Choi transfer data.  They are censored.  The physical Markov trajectory may
continue, but the exact Choi observable is intentionally set to NaN after the
first numerical-instability failure.

\section{Confidence}

Confidence is high that the regularized rank-one Choi contraction is
mathematically and programmatically represented consistently.  Confidence is
moderate that the early-cycle \(DW=1\) slab transfer evidence indicates a
near-marginal transfer channel relative to the homogeneous control.  Confidence
is low for any long-time or thermodynamic conclusion from this specific
campaign, because all locally available samples are censored before the final
cycle and the evidence inspected here is aggregate rather than raw spectral
data.

\end{document}
